\section{Distinction provenance in action}
\label{sec:demonstrations}

\subsection{Four located mechanisms}

The model and evidence package exhibit four distinct outcomes under one common audit.

\begin{table}[H]
\centering
\caption{Distinction provenance across the audited model.}
\label{tab:provenance}
\small
\begin{tabularx}{\linewidth}{@{}p{0.20\linewidth}p{0.25\linewidth}X@{}}
\toprule
Mechanism & Diagnostic object & Located result\\
\midrule
Transition folding &
Phase preimages of one target &
Three complete pre-states share one future; the loss occurs in the phase advance.\\
Boundary clipping &
Two in-domain overshoots &
States at distance \(0.2\) become identical under the clamp; the loss occurs at the wall operation.\\
Observation fibre &
Anchor-labelled phrase &
160 states produce 103 phrases; the loss occurs in the readout.\\
Finite access &
Endpoint inverse in float64 &
Recovery error grows with horizon on unique certified branches; the loss is practical access at a chosen tolerance.\\
\bottomrule
\end{tabularx}
\end{table}

These outcomes are not competing explanations of the same event. The core transition and an optional boundary or controller act in sequence; description and numerical storage are separate representations of the retained complete state:
\begin{equation}
\boxed{
\begin{aligned}
x&\xrightarrow{\ \Psi_0\ }u\xrightarrow{\ \mathcal B\ }x^+,
&&\text{core advance and optional boundary},\\
x^+&\xrightarrow{\ \Obs\ }y,
&&\text{description branch},\\
x^+&\xrightarrow{\ \mathcal N\ }\widehat{x}^{+},
&&\text{stored numerical-state branch}.
\end{aligned}
}
\label{eq:pipeline}
\end{equation}
Here \(\mathcal B\) may be the identity when no boundary or controller is present. The provenance question asks for the first map at which a chosen distinction is merged or becomes inaccessible.

\subsection{Boundary clipping as an exact operation}

\begin{proposition}[Wall-clamp collision]
\label{prop:wall}
Let one force-free coordinate have wall \(w\), equal velocities \(v>0\), and two in-domain positions \(r_1\neq r_2<w\) satisfying
\[
r_1+v\dt>w,
\qquad
r_2+v\dt>w.
\]
If the boundary rule replaces every overshoot by \(w\) and multiplies both velocities by the same bounce factor, the two states have one image.
\end{proposition}

\begin{proof}
The clamp sends both advanced positions to \(w\); equal velocities receive the same update; all other coordinates are equal.
\end{proof}

The computational construction uses \(w=3.6\), pre-state distance \(0.2\), and produces final distance zero. The proposition gives a precise engineering lesson: an invertible boundary must retain overshoot, for example by periodic wrapping on a torus or exact specular reflection.

\subsection{Relational phase transport}

The two-clock law \cref{eq:two-clock} also yields an exact retained relation. When coupling is switched off after step \(n_0\),
\begin{equation}
\delta_n
=
\delta_{n_0}
+(n-n_0)\Delta\omega\,\dt
\pmod{2\pi}.
\label{eq:decoupled-phase}
\end{equation}
Equivalently, the co-rotating coordinate
\[
I_n:=\delta_n-n\Delta\omega\,\dt
\]
is conserved. Equal frequencies reduce this to conservation of the raw phase difference.

\begin{figure}[H]
\centering
\begin{subfigure}[t]{0.49\linewidth}
\centering
\begin{tikzpicture}
\begin{axis}[
  width=\linewidth,height=0.69\linewidth,
  xlabel={uncoupled steps},
  ylabel={\(\abs{\delta}\) (radians)},
  xmin=0,xmax=305,ymin=0,ymax=3.1,
  grid=major,major grid style={draw=firule},
  tick label style={font=\scriptsize},
  label style={font=\footnotesize},
  legend style={draw=none,fill=none,font=\scriptsize,
                at={(0.03,0.97)},anchor=north west}
]
\addplot[fiblue,line width=1.2pt,mark=*]
  table[x=t,y=equal_delta_rad,col sep=comma]
  {evidence/clock_memory.csv};
\addlegendentry{\(\omega_1=\omega_2\)}
\addplot[fiorange,line width=1.2pt,mark=square*]
  table[x=t,y=unequal_delta_rad,col sep=comma]
  {evidence/clock_memory.csv};
\addlegendentry{\(\omega_1\neq\omega_2\)}
\end{axis}
\end{tikzpicture}
\caption{Relative phase.}
\end{subfigure}
\hfill
\begin{subfigure}[t]{0.49\linewidth}
\centering
\begin{tikzpicture}
\begin{axis}[
  width=\linewidth,height=0.69\linewidth,
  xlabel={uncoupled steps},
  ylabel={two-clock order parameter \(R\)},
  xmin=0,xmax=305,ymin=0.2,ymax=1.03,
  grid=major,major grid style={draw=firule},
  tick label style={font=\scriptsize},
  label style={font=\footnotesize},
  legend style={draw=none,fill=none,font=\scriptsize,
                at={(0.03,0.05)},anchor=south west}
]
\addplot[fiblue,line width=1.2pt,mark=*]
  table[x=t,y=equal_R,col sep=comma]
  {evidence/clock_memory.csv};
\addlegendentry{\(\omega_1=\omega_2\)}
\addplot[fiorange,line width=1.2pt,mark=square*]
  table[x=t,y=unequal_R,col sep=comma]
  {evidence/clock_memory.csv};
\addlegendentry{\(\omega_1\neq\omega_2\)}
\end{axis}
\end{tikzpicture}
\caption{Coherence.}
\end{subfigure}
\caption{After four coupling steps, \(\beta\) is set exactly to zero. Equal-frequency phase difference is retained; unequal-frequency drift follows \cref{eq:decoupled-phase} to floating-point precision.}
\label{fig:phase-transport}
\figsource{\sourcecode{evidence/thesis\_verification.json}, \texttt{decoupled\_clock\_memory\_and\_readout}.}
\end{figure}

\subsection{Controlled local response}

Resetting one phase is an intervention. Clone one phase state into a reset arm and a no-reset control, alter only the separation, and compare the partner after one step. The exact difference is
\begin{equation}
\Delta_{\rm reset}\phi_2^+
=
\frac{\beta\dt}{\widetilde D}
\left[
\sin(\phi_{1,\rm reset}-\phi_2)
-
\sin(\phi_{1,\rm control}-\phi_2)
\right].
\label{eq:reset-response}
\end{equation}
Measured effects are \(0.32266\) radians at distance \(0.3\) and \(0.004047\) radians at distance \(24\), matching the softened \(1/\widetilde D\) ratio. The partner does not change in the same instant; the next advance transmits the controlled response. This supplies a clean example of relational memory, local intervention, and distance-weighted dynamics within one exact phase law.

\subsection{What the demonstrations establish}

The multiroot, wall, phrase, and round-trip experiments form a single positive result: the framework identifies \emph{which operation} owns each loss of distinguishability. This is stronger than labeling every ambiguity ``entropy'' or every failed reconstruction ``irreversibility.'' It produces a typed result that can be carried into other dynamical, observational, and computational systems.
